\section{Experimental setting}\label{experimental-setting}

\textbf{Note:} The following describes a proposed simulation architecture for testing the brane model. Implementation is ongoing and numerical results will be presented in future work. This section serves as a design specification outlining the computational framework, numerical methods, and falsifiable predictions that the model generates.

\subsection{Implementation Philosophy: Emergent vs. Explicit}\label{implementation-philosophy}

\textbf{Core principle:} The simulation integrates \emph{only} the microscopic
brane degrees of freedom (node positions and velocities) evolving according to
the discretized equations of motion derived from the action
\eqref{eq:brane-action}. All emergent quantities
— effective metric, electromagnetic fields, particle identities, measurement
outcomes, etc. — are computed diagnostically from $\mathbf X(x,t)$ and
$\dot{\mathbf X}(x,t)$ and are \emph{never} fed back into the equations of motion
as independent sources, forces, or constraints.

\paragraph{What is implemented (explicit in code):}
\begin{itemize}
\item The 4D embedding map $\mathbf{X}(x,t) \in \mathbb{R}^4$ discretized as node positions $\mathbf{R}_p(t)$
\item Elastic energy functionals: stretching $U_\text{str}$ and bending $U_\text{bend}$
\item Time evolution via velocity-Verlet integration: $\mathbf{F}_p = -\partial U/\partial \mathbf{R}_p$, $m_p \ddot{\mathbf{R}}_p = \mathbf{F}_p$
\item Boundary conditions (periodic, free, or clamped)
\item Initial conditions: Gaussian pulses, toroidal configurations, amplitude patterns
\end{itemize}

\paragraph{What must emerge (detected via measurement, \emph{not} coded):}
\begin{itemize}
\item \textbf{Lorentz invariance:} Isotropic dispersion $\omega(|\mathbf{k}|)$ measured from wave propagation tests (Section~\ref{calibration-dispersion-isotropy})
\item \textbf{Gravity:} Induced metric $g_{ij} = \partial_i \mathbf{X} \cdot \partial_j \mathbf{X}$ computed post-hoc from node positions; geodesic deflection ray-traced through $g$
\item \textbf{Electromagnetism:} Coulomb-like $1/r^2$ field measured from time-averaged amplitude $\bar{X}^4(\mathbf{x})$ around stable toroidal solitons
\item \textbf{Nonlinear threshold:} Emergent energy density threshold for soliton nucleation; \emph{not} imposed as a hard cutoff on amplitude or strain
\item \textbf{Particle stability:} Toroidal solitons that self-confine via nonlinear elasticity; no explicit "particle glue" or confinement potential added
\item \textbf{Spinor behavior:} $4\pi$ holonomy measured from phase transport around rotated configurations
\end{itemize}

\paragraph{No back-reaction from emergent fields.}
Even when we visualise or analyse effective quantities such as $g_{ij}$,
Coulomb-like fields, or reconstructed wavefunctions, these objects remain purely
diagnostic. In the implementation they are read out from the brane state for
measurement and comparison with continuum theories; they never act back on the
brane as additional forces or constraints. This ensures that any agreement with
relativity, electromagnetism or quantum statistics is genuinely \emph{emergent}
from the substrate dynamics.

\paragraph{Photon field initialization (clarification):}
When we refer to "implementing a photon field component" for the electron model, this means \emph{initializing} the brane with a photon-like wave pattern confined in a toroidal structure—a specific choice of $\mathbf{R}_p(t=0)$ and $\mathbf{v}_p(t=0)$. The subsequent time evolution is governed \emph{solely} by the elastic forces $\mathbf{F}_p = -\partial U/\partial \mathbf{R}_p$. There is \textbf{no electromagnetic force} $\mathbf{F}_\text{EM}$ added to the equations of motion; any EM-like behavior (field gradients, Coulomb scaling) must appear naturally in the resulting amplitude distribution.

\paragraph{Dirac-consistent rest-frame initial data.}

For the numerical experiments in this work we primarily consider a
rest-frame electron with a definite spin orientation. In the brane
picture this is realized as a toroidal soliton whose centerline lies
in a fixed lateral plane, with the torus axis aligned with a chosen
quantization axis (taken to be the $X^3$-direction for definiteness).
The internal amplitude mode $\xi(t,z,x,y)$ in the corresponding
world-tube is initialized as a Compton-frequency oscillation that
circulates at the microscopic brane wave speed $c$ along the tube, as
described in Section~\ref{subsec:tubular-electron}.

In the present numerical experiments we adopt the tubular ansatz
of Eq.~\eqref{eq:xi-electron-init}, implementing a W\&vdM-inspired
intertwined double loop. The cross-section profile
$f(\rho,\theta;z)$ from Eq.~\eqref{eq:double-tunnel-profile} describes
two narrow ridges of high amplitude (our colloquial ``energy tunnels'')
at angular positions $\theta = \alpha(z)$ and
$\theta = \alpha(z)+\pi$ inside the toroidal region of the brane, where
$\alpha(z) = \phi(z) = z/R$ controls the twist of the ridge pair as one
moves along the torus. With twist winding $\ell=1$, each ridge executes
a single rotation as the torus is traversed once, producing an
intertwined $(2,1)$-type torus-knot pattern in which the two ridges
spiral around each other. The longitudinal phase factor with
winding number $m=2$ encodes that the internal Compton mode winds twice
around the torus per revolution. This geometry closely mirrors the
W\&vdM double-loop picture while remaining fully compatible with the
brane's elastic dynamics.

The Dirac equation enters at this stage as an effective constraint on
the coarse-grained observables of the soliton rather than as a
microscopic wave equation. In particular, for a rest-frame spin-up
electron the Dirac bilinears imply a four-current with vanishing net
spatial component and a spin three-vector of magnitude $\hbar/2$
aligned with the quantization axis, cf.\
Eq.~\eqref{eq:dirac-rest-bilinears}. We impose the corresponding
conditions on the toroidal ansatz by choosing the axis of the torus,
the direction of the internal circulation of $\xi$, and a small
tangential lateral velocity field in the tube such that the resulting
coarse-grained momentum and spin of the soliton approximate these
Dirac values. The subsequent evolution is always governed by the full
nonlinear brane equations; the Dirac constraints are used only to fix
and assess the macroscopic properties (charge, spin, effective
four-current) of the initial electron state.

\paragraph{Experimental workflow:}
\begin{enumerate}
\item Set initial conditions $\{\mathbf{R}_p(0), \mathbf{v}_p(0)\}$ representing a physical scenario (e.g., localized pulse, toroidal electron, point mass)
\item Evolve forward using only elastic forces from $U_\text{str} + U_\text{bend}$
\item \textbf{Measure} emergent quantities: dispersion $\omega(\mathbf{k})$, induced metric $g_{ij}$, energy localization, field profiles $\bar{X}^4(\mathbf{x})$, geodesic deflection angles, etc.
\item Compare measured behavior against known physics (Lorentz cones, $1/r$ potentials, $1/r^2$ fields) to assess whether the substrate dynamics reproduce standard results
\end{enumerate}

This approach ensures that the simulation tests the \emph{sufficiency} of geometric elasticity to reproduce observed physics, rather than presupposing the structure of electromagnetism, gravity, or quantum thresholds.

\paragraph{No feedback from emergent fields (architectural constraint).}
For clarity, we emphasize a strict one-way flow of information in the
implementation. The microscopic solver evolves only the brane embedding
$\mathbf{X}(x,t)$ under the elastic energy functionals described in
Sec.~\ref{discrete-energies-and-forces}, using a fixed integration scheme and
boundary conditions. All effective fields --- electromagnetic
potentials and fields, induced metrics, gravitational potentials, and
Lorentz-force-like terms --- are defined purely as \emph{diagnostic
functionals} of $\mathbf{X}$ and its derivatives, evaluated on
snapshots of the simulation. They are never fed back into the
microscopic equations of motion as additional forces, constraints, or
dimension-specific damping terms. In particular, the fourth (normal)
component $X^4$ enters the dynamics only through the mechanical
Lagrangian; any identification of $X^4$ with an effective charge or
scalar potential is an interpretational step applied \emph{after} the
fact. This architectural separation ensures that the numerical
experiments genuinely test whether the simple brane dynamics are
\emph{sufficient} to reproduce the observed structures of relativity,
electromagnetism, and gravity.

\subsection{Measurement procedures: spectral band isolation and Berry reconstruction}
\label{measurement-procedures}

A key methodological change in the new layered framework is that the effective
electromagnetic sector is not reconstructed from a direct coarse-graining of
the raw normal displacement $X^4$, but from the geometric phase structure of
\emph{band-isolated} substrate modes. Concretely, we reconstruct an effective
four-potential from the pulled-back Berry connection introduced in
Section~\ref{sec:berry-pullback-em}.

\subsubsection{Data products from the substrate solver}

At each time step the solver provides node positions and velocities,
$\mathbf R_p(t)\approx \mathbf X(x_p,t)\in\mathbb R^4$ and
$\dot{\mathbf R}_p(t)\approx \partial_t\mathbf X(x_p,t)$.
All higher-layer fields are computed diagnostically from these arrays and their
finite differences.

\subsubsection{Local space--time spectral analysis}

To obtain a meaningful complex envelope from a real substrate, we apply a
\emph{windowed} space--time Fourier analysis (short-time Fourier transform / local
FFT). On a chosen spatial patch and within a time window of length $\Delta t$ we
compute $\hat{\xi}^A(\mathbf k,\omega)$ as in Eq.~\eqref{eq:xi-fourier}, using
tapered windows to control leakage.

The output of this step is a complex spectrum that resolves:
(i) spatial wavevector $\mathbf k$,
(ii) temporal frequency $\omega$, and
(iii) a four-component complex mode amplitude.

\subsubsection{Band identification and projection (Route 2 in practice)}

For the same parameter settings we construct the linearized dynamical matrix
$D(\mathbf k)$ of the discrete substrate (details depend on the discrete energy
model in Section~\ref{discrete-energies-and-forces}) and compute its
eigenpairs $(\omega_n(\mathbf k),u_n(\mathbf k))$.

Given a target band index $n$ (``photon-like'' band), we project the measured
spectrum onto the polarization eigenvector:
\begin{equation}
\hat{\psi}_n(\mathbf k,\omega)
=\langle u_n(\mathbf k),\hat{\xi}(\mathbf k,\omega)\rangle,
\label{eq:band-projection}
\end{equation}
and reconstruct the band-limited complex envelope $\psi_n(x,t)$ by inverse
transform restricted to that band. This yields a complex field with a single
local $U(1)$ phase, but we emphasize again: this phase is not yet the Berry
phase; Berry structure is computed from the geometry of $u_n(\mathbf k)$.

\subsubsection{Numerical Berry connection and curvature}

We compute the Berry connection and curvature for the isolated band using a
gauge-robust discrete formulation on a $\mathbf k$-grid (link-variable method).
This avoids numerical instability from explicit derivatives
$\partial_{k_i}u_n$ and ensures invariance under $u_n\mapsto e^{i\chi}u_n$.

The result is $\mathcal A_i(\mathbf k)$ and $\mathcal F_{ij}(\mathbf k)$ as in
Eqs.~\eqref{eq:berry-connection}--\eqref{eq:berry-curvature}.

\subsubsection{Reconstruction of an effective electromagnetic field}

From the band envelope $\psi_n(x,t)$ we estimate a slowly varying local
wavevector field $\mathbf k(x,t)$ (either via local FFT peak tracking or via the
phase-gradient estimator $\mathbf k\approx \nabla \arg\psi_n$ in regions of high
signal-to-noise). We then compute the pulled-back Berry one-form
$\mathcal A_\mu(x,t)$ from Eq.~\eqref{eq:berry-pullback} and define
$A_\mu^{\mathrm{EM}}$ and $F_{\mu\nu}^{\mathrm{EM}}$ via
Eqs.~\eqref{eq:em-potential-from-berry}--\eqref{eq:em-tensor-from-berry}.

We extract the usual electric and magnetic fields as
\begin{equation}
E_i = F_{0i}^{\mathrm{EM}},
\qquad
B_i = \frac{1}{2}\epsilon_{ijk}F_{jk}^{\mathrm{EM}}.
\end{equation}
The homogeneous Maxwell equations hold identically by construction (Bianchi
identity). The inhomogeneous equations become part of the numerical test plan:
we measure whether effective sources localize at defects/solitons and whether
wave dynamics reproduces the expected vacuum behaviour away from defects.

\subsubsection{Gravitational diagnostics (unchanged principle)}

Gravitational effects are still reconstructed from the induced metric
$g_{ij}=\partial_i\mathbf X\cdot\partial_j\mathbf X$ and ray tracing in the
effective interval $ds^2=-c^2dt^2+g_{ij}(x,t)\,dx^i dx^j$ (see
Section~\ref{gravity-as-effective-metric}). This part of the pipeline remains
independent of the electromagnetic reconstruction above.

\subsection{Discrete Simulation Framework}\label{discrete-simulation-framework}

We discretize the material domain $\Omega$ by nodes $x_p\in\Omega$ and store
\begin{equation}\label{eq:discrete-embedding}
  \mathbf R_p(t) \approx \mathbf X(x_p,t)\in\mathbb R^4,
\end{equation}
so that the discrete configuration $\{\mathbf R_p\}$ is a sampling of the continuum embedding
$\mathbf X(x,t)$ introduced in Section~\ref{sec:conceptual-model}.
The fourth component of $\mathbf R_p$ encodes the out-of-brane displacement; there is no separate amplitude field.

The choice of continuum parameters $(\rho_m,T,\kappa)$
is constrained by the parameter-matching conditions of
Section~\ref{sec:param-matching}; in particular, the amplitude scale calibration
of Section~\ref{amplitude-scale-calibration} implements the matching to
$(m_e,c,\hbar,e)$ at the level of the discrete lattice.


\subsubsection{Lattice and neighborhoods}\label{lattice-and-neighborhoods}

\begin{itemize}
\item Use an \textbf{FCC lattice} in material space (spacing $h$), giving 12 equal-length nearest neighbors $\mathcal N(p)$ for good isotropy.
\item If FCC is impractical, a cubic 26-stencil with \textbf{isotropic weights} (4th-order Laplacian fit) can be used; document the weights.
\end{itemize}

\subsubsection{Discrete energies and forces}\label{discrete-energies-and-forces}

We approximate the continuum energy in terms of \textbf{edge stretches} (intrinsic strain) and \textbf{bending} between neighboring edge directions.

\begin{itemize}
\item \textbf{Stretching energy (Hooke-like):}
  \[
  U_\text{str}=\sum_{(p,q)\in\mathcal E}
  \phi\left(\frac{|\mathbf R_q-\mathbf R_p|-h}{h}\right),
  \]
  with a simple quadratic potential
  \[
  \phi(\epsilon)=\frac{1}{2}k\,\epsilon^2,
  \]
  so that the discrete stretching law reduces locally to a
  Hooke-like response. In the present work we do not include
  any additional material saturation; all nonlinear behavior
  arises from the geometric dependence of distances and curvature
  on the brane embedding.

\paragraph{Nonlinearity mode (implementation note).}

In code we distinguish multiple nonlinearity modes via a
\texttt{BraneNonlinearityMode} enumeration:
\begin{itemize}
  \item \texttt{NONE}: Linearized small-displacement approximation;
        edge stiffness computed using $\epsilon \approx
        (\partial_i\xi_j)/h$ rather than the full 4D Euclidean
        distance.
  \item \texttt{PURE\_GEOMETRY}: Default. Each edge stiffness is
        computed from the exact 4D Euclidean distance
        $\ell=|\mathbf R_q-\mathbf R_p|$, so lateral contraction from
        normal displacements arises purely from metric changes. This is
        the primary mode for all emergent-field experiments and
        corresponds exactly to the theoretical model developed in
        Section~\ref{conceptual-model}.
  \item \texttt{SCHWINGER\_LIMIT}, \texttt{HUYGENS\_ELASTIC},
        \texttt{HIGGS\_MEXICAN\_HAT}, \texttt{CUSTOM}: \emph{Purely
        exploratory} extensions for future robustness tests. These modes
        introduce additional material-nonlinearity potentials beyond the
        geometric coupling and are \textbf{not part of the core theoretical
        framework}. They are not used in any results presented in this paper.
\end{itemize}
Unless explicitly stated otherwise, all simulations and numerical experiments
reported in this work use \texttt{PURE\_GEOMETRY}.

  The derivative (useful for force implementation) is
  \[
  \phi'(\epsilon) = k\epsilon.
  \]

\item \textbf{Bending/curvature energy (optional but stabilizing):}
  For each node $p$, penalize changes of edge directions:
  \[
  U_\text{bend}=\frac{\kappa}{2}\sum_{p}\sum_{q<r\in\mathcal N(p)}
  \big(1-\hat e_{pq}\cdot\hat e_{pr}\big),
  \quad \hat e_{pq}=\frac{\mathbf R_q-\mathbf R_p}{|\mathbf R_q-\mathbf R_p|}.
  \]
  For small angles between neighboring edges, $1-\hat e_{pq}\cdot\hat e_{pr}\approx \tfrac12\theta^2$.
  Thus $U_\text{bend}$ acts as a discrete approximation to the squared mean curvature $|b|^2$ of the
  embedding, consistent with the continuum brane picture. This helps maintain coherent toroidal shapes.

\item \textbf{Total energy:} $U=U_\text{str}+U_\text{bend}$.
\end{itemize}

Forces are $\mathbf F_p=-\partial U/\partial \mathbf R_p$. Because $|\mathbf R_q-\mathbf R_p|$ is a \textbf{4D norm}, any deformation in the 4th direction automatically couples to in-brane edge lengths, producing the desired \textbf{geometric lateral contraction}.

\paragraph{Discrete-to-continuum parameter mapping}\label{discrete-continuum-mapping}

The continuum parameters introduced in Section~\ref{conceptual-model} relate to the discrete simulation as follows. Each lattice node represents a material volume element of size $h^3$, where $h$ is the lattice spacing (edge length in the reference FCC configuration). The node mass is $m_p = \rho_m h^3$, with $\rho_m$ the continuum mass density of the brane substrate. The effective \textbf{isotropic tension} $T$ in the continuum is encoded in the discrete spring stiffness parameter $k$ appearing in $\phi(\epsilon)$; in the small-strain limit, dimensional analysis and dispersion measurements yield $T \sim k/h$. The \textbf{bending stiffness} $\kappa$ in the continuum curvature energy $\tfrac{\kappa}{2}|b|^2$ maps to the discrete parameter $\kappa$ in $U_\text{bend}$, with units adjusted by powers of $h$ to maintain dimensional consistency. Finally, the shear modulus $\mu$ in the continuum description can be introduced via additional discrete angular or shear springs if needed, though for simplicity we often set $\mu\to 0$ in the isotropic limit. These mappings are validated by measuring the wave speed $c = \sqrt{T/\rho_m}$ from numerical dispersion tests and verifying consistency with the continuum prediction (see Section~\ref{calibration-dispersion-isotropy}).

\subsubsection{Amplitude scale calibration}\label{amplitude-scale-calibration}

To instantiate the brane in physical units we must fix not only the wave speed
$c$ but also the absolute scale of normal displacements $X^4$, which has the
same dimension of length as the lateral coordinates. Two complementary routes
are used to estimate the amplitude scale $A$ of microscopic oscillations.

\paragraph{Route (i): Compton-cell mass calibration.}
We first tie the bulk mechanical parameters to the electron rest mass and the
reduced Compton wavelength
$\lambda_C = \hbar/(m_e c)$. In the continuum, a small-amplitude transverse wave
with typical amplitude $A$, angular frequency $\omega$ and effective volume
$V_{\text{eff}}$ has elastic energy
\begin{equation}
  E_{\text{wave}} \;\sim\; \frac{1}{2}\,\rho\,A^2\,\omega^2\,V_{\text{eff}},
  \label{eq:wave-energy-amplitude}
\end{equation}
where $\rho$ is the brane mass density. As a simple calibration ansatz we assume
that a Compton-scale ``cell'' of the brane carries roughly one electron mass,
\begin{equation}
  \rho\,\lambda_C^3 \;\approx\; m_e,
\end{equation}
which fixes $\rho$ and, via $c^2 = K/\rho$, the bulk modulus $K$. For a mode
with $\omega \approx 2\pi c/\lambda_C$ localized in a volume
$V_{\text{eff}}\sim\lambda_C^3$, imposing that this configuration represents
a single quantum with energy $E_{\text{wave}}\approx\hbar\omega$ yields an
amplitude scale $A=\mathcal{O}(\lambda_C)$; a simple estimate gives
$A \approx \lambda_C/\sqrt{\pi}$ up to order-unity geometric factors. This
choice sets the overall substrate mechanical scale (via $\rho$ and $K$) in terms
of $(m_e,c,\hbar)$.

\paragraph{Route (ii): charge calibration (deferred in the new framework).}
In the previous draft, the charge magnitude was tied directly to a coarse-grained
normal-displacement monopole. In the present layered framework, the effective
electromagnetic sector is reconstructed from Berry geometry of band-isolated modes.
This shifts the charge calibration problem to fixing the conversion factor
$\hbar/q$ in Eq.~\eqref{eq:em-potential-from-berry} and identifying how defects
realize effective sources in the inhomogeneous Maxwell sector. We therefore
defer a first-principles charge matching to the dedicated numerical tests of
Section~\ref{falsifiable-criteria}.

\paragraph{Practical simulation choice.}
In the discrete code we select a lattice spacing $h$ and identify each node with
a material volume $h^3$ and mass $m_p = \rho h^3$. Given a target amplitude
scale $A$ from either route, a ``one-quantum'' initialization of a photon-like
packet or a toroidal electron is implemented by choosing discrete displacements
$[R_p]_4$ such that the discrete elastic energy
$\tfrac{1}{2}\sum_p m_p |\dot{R}_{p,4}|^2 + U_{\text{elastic}}(\{\mathbf R_p\})$
matches $\hbar\omega$ (for a photon) or $m_e c^2$ (for an electron) to within a
specified tolerance. Both calibration routes therefore determine not only the
mechanical parameters $(\rho,K)$ but also the physically meaningful range of
normal displacements $X^4$ to be used when setting up initial conditions in
numerical experiments.

\paragraph{SI units to dimensionless lattice units.}
The discrete simulation operates in dimensionless units internally, with all
physical quantities scaled by characteristic Compton-scale parameters. The
mapping from SI units to simulation units proceeds as follows:
\begin{itemize}
  \item \textbf{Length:} $[x]_{\text{sim}} = x/\lambda_C$, so that one lattice
    spacing $h \sim 0.1$--$1.0$ in simulation units corresponds to
    $h_{\text{SI}} \sim 0.1\lambda_C$--$1.0\lambda_C$ physically
    ($\lambda_C \approx 3.86 \times 10^{-13}$\,m for the electron).
  \item \textbf{Time:} $[t]_{\text{sim}} = t\,\omega_C = t c/\lambda_C$, where
    $\omega_C = m_e c^2/\hbar$ is the Compton angular frequency. One Compton
    period $2\pi/\omega_C \approx 1.29 \times 10^{-21}$\,s maps to
    $[t]_{\text{sim}} = 2\pi$.
  \item \textbf{Mass density:} $[\rho]_{\text{sim}} = \rho\,\lambda_C^3/m_e$,
    normalized such that a Compton-scale cell carries one electron mass. The
    continuum calibration $\rho\,\lambda_C^3 \approx m_e$ corresponds to
    $[\rho]_{\text{sim}} \approx 1$.
  \item \textbf{Energy:} $[E]_{\text{sim}} = E/(m_e c^2)$, so that a rest
    electron carries $[E]_{\text{sim}} = 1$ and a photon with frequency
    $\omega_C$ carries $[E]_{\text{sim}} = 1$. In SI units, $m_e c^2 \approx
    8.19 \times 10^{-14}$\,J.
\end{itemize}
With these scalings, the wave speed $c$ becomes $[c]_{\text{sim}} = 1$
(measured in lattice spacings per Compton period), and $\hbar$ is implicitly
set to unity in natural units. All forces, energies, and velocities in the code
are computed in these dimensionless units; conversion back to SI units for
comparison with experiment requires multiplying by the appropriate power of
$(\lambda_C, \omega_C^{-1}, m_e)$. This Compton-based scaling makes the
simulation independent of the absolute physical scale while preserving all
relevant dimensionless ratios.

\subsubsection{Time integration and stability}\label{time-integration-and-stability}

\begin{itemize}
\item \textbf{Integrator:} velocity-Verlet (leapfrog), symplectic \& time-reversible:
  \begin{align*}
  \mathbf v^{n+1/2}&=\mathbf v^n+\tfrac{\Delta t}{2}\mathbf a^n,\\
  \mathbf R^{n+1}&=\mathbf R^n+\Delta t\mathbf v^{n+1/2},\\
  \mathbf v^{n+1}&=\mathbf v^{n+1/2}+\tfrac{\Delta t}{2}\mathbf a^{n+1}.
  \end{align*}
\item \textbf{CFL guideline (linear regime):} $\Delta t<\eta h/c$ with empirical $\eta\approx0.33$ for FCC. Here $c$ is measured from dispersion tests (Section~\ref{calibration-dispersion-isotropy}) rather than assumed from continuum $c=\sqrt{T/\rho}$. Document the chosen $\eta$ and measured $c$.
\item \textbf{Energy tracking:} $E_\text{tot}=K+U$ with $K=\tfrac12\sum_p \rho h^3|\dot{\mathbf R}_p|^2$.
  Here $\rho$ is the continuum mass density and $h^3$ is the material volume represented by a node,
  so that each node effectively carries mass $m_p=\rho h^3$.
  Report drift (target $<10^{-5}$ over long runs, periodic BCs).
\item \textbf{Implementation note:} No explicit amplitude or strain clamps are applied during time integration. All nonlinear confinement effects arise from the geometric dependence of edge lengths on the 4D embedding. The integrator advances $\mathbf R_p$ and $\mathbf v_p$ based solely on forces $\mathbf F_p = -\partial U/\partial \mathbf R_p$ computed from $U_\text{str} + U_\text{bend}$, ensuring that any emergent localization or threshold behavior arises from the elastic dynamics rather than from artificial hard cutoffs.
\end{itemize}

\subsubsection{Boundary conditions}\label{boundary-conditions}

In the discrete simulations we must close the brane patch by choosing
boundary conditions that specify how the medium interacts with the
“outside world”. In practice we use:

\begin{itemize}
\item \textbf{Periodic} (default numerics): the simulation domain is
  wrapped onto a torus so that waves leaving one side re-enter on the
  opposite side. This mimics a locally homogeneous, effectively infinite
  brane patch.
\item \textbf{Free} ($\partial_n \mathbf X$ traction-free in the in-brane
  sense): the brane is not anchored at the boundary, allowing transverse
  motion without external forces; this approximates an isolated finite
  membrane embedded in a larger medium.
\item \textbf{Clamped} ($\mathbf X=\mathbf X_0$ on $\partial\Omega$):
  the boundary is rigidly fixed, representing attachment to an external
  structure and suppressing fluctuations at the edge.
\end{itemize}

\subsubsection{Induced metric and geodesics (from data)}\label{induced-metric-and-geodesics}

Given positions $\mathbf R_p$, estimate discrete tangents $\partial_i\mathbf X$ via central differences in material space, build $g_{ij}=\partial_i\mathbf X\cdot\partial_j\mathbf X$, and ray-trace probe signals using $ds^2=-c^2dt^2+g_{ij}dx^i dx^j$. This quantifies "gravitational" deflection due to 4D bulges.

\subsubsection{Calibration: dispersion \& isotropy (emergent relativity test)}\label{calibration-dispersion-isotropy}

Initialize small-amplitude plane waves $\mathbf R_p=\mathbf R^0_p+\boldsymbol\xi\cos(\mathbf k\cdot x_p)$ in three non-collinear $\hat k$ and \emph{numerically measure} $\omega(\mathbf k)$ from the time evolution. Verify that the measured dispersion follows
\[
\omega \approx c|\mathbf k|,\qquad
\text{relative error } \le 1\% \text{ up to } kh\approx0.5.
\]
This demonstrates isotropic cones for internal observers in the linear regime.

\paragraph{Diagnostics and observables.}
In numerical experiments we monitor (i) wave propagation and dispersion of
photon-like excitations, (ii) localization and stability of soliton-like
deformations, (iii) effective metric properties inferred from signal
propagation, and (iv) coarse–grained energy densities that play the role of
effective mass sources. All of these are computed diagnostically from
$\mathbf X(x,t)$ and its time derivative; no emergent field is fed back into
the microscopic dynamics.

\paragraph{1D photon test with lateral distortion.}
As a minimal check of the geometric nonlinearity in a realistic regime, we
implemented a one–dimensional spring chain with Compton–calibrated lattice
spacing, mass density and tension (fixing the wave speed to $c$). A
Gaussian wave packet in the normal direction $X^4$ propagates across the
chain with essentially unchanged shape and speed, reproducing the linear
photon picture at realistic scales. Because the springs are slightly
pre-tensioned ($L_0 < h$), the same run also exhibits a small, localized
co–moving distortion in the lateral separation of nodes. This lateral
``bulge'' is confined to the region of largest normal gradients and remains
bounded over the simulation time, illustrating how the pre-tensioned
Hookean brane generates a mild coupling between normal and lateral
coordinates without destroying photon propagation.

\subsubsection{Test plan and falsifiable criteria}
\label{falsifiable-criteria}

All tests are formulated as measurements on simulation outputs. No emergent
quantity is allowed to feed back into the substrate solver.

\begin{enumerate}
\item \textbf{Relativistic cones (substrate linear regime):}
Measure $\omega(\mathbf k)$ for small-amplitude plane waves in multiple
directions and verify isotropy and approximately linear dispersion
$\omega\approx c|\mathbf k|$ up to a documented cutoff.

\item \textbf{Band isolation sanity:}
Inject a known single-mode excitation and verify that projection
Eq.~\eqref{eq:band-projection} recovers the correct band amplitude while
suppressing leakage into other bands below a set tolerance.

\item \textbf{Gauge robustness of Berry numerics:}
Verify invariance of $\mathcal F_{ij}(\mathbf k)$ under random phase
redefinitions $u_n(\mathbf k)\mapsto e^{i\chi(\mathbf k)}u_n(\mathbf k)$.
Optionally, test quantized invariants (e.g. Chern number) on closed $\mathbf k$-surfaces
in controlled toy configurations.

\item \textbf{Berry $\rightarrow$ EM (homogeneous Maxwell laws):}
Compute $F_{\mu\nu}^{\mathrm{EM}}$ from the pulled-back Berry connection and
numerically verify $\nabla\cdot\mathbf B\approx 0$ and
$\partial_t\mathbf B+\nabla\times\mathbf E\approx 0$ away from defects, with
errors decreasing under refinement.

\item \textbf{Defect localization of effective sources:}
In runs that form a stable localized soliton (electron candidate), test whether
violations of the inhomogeneous Maxwell equations (interpreted as effective
sources) localize at the defect core and vanish in the far field.

\item \textbf{Spinorial holonomy ($4\pi$ behaviour):}
Transport the reconstructed internal frame around a closed loop enclosing the
defect and measure the holonomy. The target signature is a $4\pi$ periodicity
consistent with an emergent spinor sector.

\item \textbf{Aharonov--Bohm type holonomy (optional, medium term):}
Construct two paths of the same band envelope enclosing a region of nontrivial
Berry curvature and test whether the relative phase shift matches the path
integral of the reconstructed effective potential \cite{AharonovBohm1959}.
\end{enumerate}

\subsection{Interpretation layers (clarity note)}\label{interpretation-layers}

\begin{itemize}
\item \textbf{Substrate:} single geometric object $\mathbf X(x,t)\in\mathbb R^4$ with dynamics from $\mathcal L(\mathbf X,\partial\mathbf X)$.
\item \textbf{Emergent effective physics:} induced metric $g_{ij}$ and wave kinematics; no extra fields postulated.
\item \textbf{Threshold phenomena:} energy density criterion $\langle\mathcal E\rangle_{V_c}V_c \ge m_e c^2$ emerges from geometric nonlinearity, not imposed as external cutoff.
\end{itemize}

\textbf{Note on implementation:} The discrete simulation framework described in Section~\ref{discrete-simulation-framework} supersedes earlier spring-lattice prototypes. Those preliminary implementations used a simpler cubic lattice with Hooke's-law springs and Euler integration; the current framework uses an FCC lattice with geometric nonlinearity and velocity-Verlet integration, providing better isotropy and numerical stability.

\paragraph{Geometric nonlinearity diagnostic tool.}
For diagnostic purposes we also provide a small Python script
(used to generate Fig.~\ref{fig:geom-nonlinearity-compton}) that
computes and visualizes $\eta(\lambda)$ for realistic electron
parameters. This tool is not part of the simulation core, but serves
to illustrate the Compton-scale onset of geometric nonlinearity in the
calibrated substrate.
